\documentclass[11pt,a4paper]{article}
% head height 14pt: the italic Th ligature of a running head stands 13.01pt
\usepackage[margin=35mm,top=32mm,bottom=45mm,headheight=14pt]{geometry}
\usepackage{amsmath,amsthm,mathtools}
\usepackage{fontspec}
\usepackage[mathsf=sym]{unicode-math}
\setmainfont{EBGaramond}[Path=../fonts/, Extension=.otf,
  UprightFont=*-Regular, ItalicFont=*-Italic,
  BoldFont=*-Bold, BoldItalicFont=*-BoldItalic,
  FontFace={sb}{n}{*-SemiBold}, Numbers=Lining,
  SmallCapsFeatures={LetterSpace=4}]
% no sans text face in the family: \textsf (crate names) is code, in the mono
\defaultfontfeatures[Iosevka]{Path=../fonts/, Extension=.ttf,
  UprightFont=*-Regular, ItalicFont=*-Italic, BoldFont=*-Bold,
  FontFace={sb}{n}{*-SemiBold}, Scale=MatchLowercase}
\setsansfont{Iosevka}
\setmonofont{Iosevka}
% stylistic set 3: the calligraphic \mathcal (the default is a roundhand
% script). This copy maps capitals shared with Greek to the Latin, so copy
% and search see Latin (fonts/README-Garamond-Math.txt)
\setmathfont{Garamond-Math-Latin.otf}[Path=../fonts/, StylisticSet=3]
% the math sans is 6 % larger-eyed than the text; scaled to the text x-height
\setmathfont{Garamond-Math-Latin.otf}[Path=../fonts/,
  range=\mathsfup/{latin,Latin,num}, Scale=0.94]
% each \mathsf is one atom, as with a classic math alphabet, so a subscripted
% identifier such as \rivk takes a further subscript
\AtBeginDocument{\let\arbsymsf\mathsf\protected\def\mathsf#1{{\arbsymsf{#1}}}}
\usepackage{polyglossia}
\setmainlanguage[variant=british,ordinalmonthday=false]{english}
\PolyglossiaSetup{english}{frenchspacing=true}
\usepackage[protrusion=true,expansion=false]{microtype}
\usepackage{booktabs,enumitem}
\usepackage{placeins}
\usepackage{tikz}
\usetikzlibrary{arrows.meta,positioning,calc,fit,backgrounds}
\definecolor{tierspend}{RGB}{200,60,60}
\definecolor{tierfull}{RGB}{210,140,40}
\definecolor{tierin}{RGB}{60,130,180}
\definecolor{tierrecv}{RGB}{70,150,90}
\usepackage[psdextra,hidelinks]{hyperref}

\emergencystretch=3em
\hbadness=10000
\hfuzz=2pt

\newtheorem{theorem}{Theorem}[section]
\newtheorem{lemma}[theorem]{Lemma}
\newtheorem{proposition}[theorem]{Proposition}
\newtheorem{corollary}[theorem]{Corollary}
\theoremstyle{definition}
\newtheorem{definition}[theorem]{Definition}
\newtheorem{construction}[theorem]{Construction}
\newtheorem{assumption}[theorem]{Assumption}
\newtheorem{remark}[theorem]{Remark}
\newtheorem{example}[theorem]{Example}

% kind-coded theorem blocks: the rule weight codes the kind, in ink, no fill:
% proved statements 1.6pt, definitions and constructions 0.4pt, assumptions
% dashed 0.8pt; remarks and examples run unframed with an italic head
\usepackage{etoolbox,framed}
\makeatletter
\providecolor{tierin}{RGB}{60,130,180}
\providecolor{tierfull}{RGB}{210,140,40}
\providecolor{tierspend}{RGB}{200,60,60}
\providecolor{muted}{HTML}{5E5E5E}
% a left rule 0.9em from the text; framed spaces it like a list. Inside,
% \topsep is cleared so the statement's own list adds nothing, and
% FrameRestore's \sloppy is dropped so statements break like prose. A frame
% split across pages counts the last line's depth, which framed's
% \boxmaxdepth=0 turns into height
\setlength{\OuterFrameSep}{7pt plus 2pt minus 2pt}
\patchcmd{\fb@put@frame}{\splittopskip\z@}{\splittopskip\z@\splitmaxdepth\z@}{}{\PatchFailed}
% the frame sits flush left, as framed would centre it
\def\arb@flushleft#1{\hbox to\hsize{#1\hfil}}
\patchcmd{\fb@putboxa}{\centerline}{\arb@flushleft}{}{\PatchFailed}
\patchcmd{\fb@put@frame}{\centerline}{\arb@flushleft}{}{\PatchFailed}
\newcommand\arb@kind[2]{%
  \BeforeBeginEnvironment{#2}{\def\FrameCommand{#1}%
    \MakeFramed{\advance\hsize-\width\let\sloppy\relax\FrameRestore\topsep\z@}}%
  \AfterEndEnvironment{#2}{\endMakeFramed}}
\newcommand\arb@dashed[1]{\setbox\z@\hbox{#1}\hbox{\vbox to\ht\z@{\cleaders
  \vbox{\hrule\@width.8pt\@height3pt\vskip2pt}\vfill}\hspace{.9em}\box\z@}}
\arb@kind{\vrule\@width.4pt\hspace{.9em}}{definition}
\arb@kind{\vrule\@width.4pt\hspace{.9em}}{construction}
\arb@kind{\arb@dashed}{assumption}
\arb@kind{\vrule\@width1.6pt\hspace{.9em}}{theorem}
\arb@kind{\vrule\@width1.6pt\hspace{.9em}}{lemma}
\arb@kind{\vrule\@width1.6pt\hspace{.9em}}{proposition}
\arb@kind{\vrule\@width1.6pt\hspace{.9em}}{corollary}
\AtBeginEnvironment{remark}{\thm@headfont{\itshape}}
\AtBeginEnvironment{example}{\thm@headfont{\itshape}}
% a proof head breaks and hyphenates as the text it opens: amsthm sets it as a
% list label, one unbreakable box
\renewenvironment{proof}[1][\proofname]{\par\pushQED{\qed}\normalfont
  \topsep6\p@\@plus6\p@\relax\trivlist\item\relax{\itshape#1\@addpunct{.}}%
  \hskip\labelsep\ignorespaces}{\popQED\endtrivlist\@endpefalse}
% 11/14.5 pt; small (captions, tables, footnotes) 0.875; both size commands
% are robust, so the patches go to their inner macros. A 10pt class (the
% figure render's standalone) keeps its sizes
\ifnum\@ptsize=1
\expandafter\patchcmd\csname normalsize \endcsname{{13.6}}{{14.5}}{}{\PatchFailed}
\expandafter\patchcmd\csname small \endcsname{\@xpt\@xiipt}{{9.625}{12.7}}{}{\PatchFailed}
\fi
% the hyperref package wraps \@footnotetext; the size is set in the wrapped
% original
\patchcmd{\H@@footnotetext}{\footnotesize}{\small}{}{\PatchFailed}
\normalsize
\setlength{\parindent}{1.2em}
% operator spacing shrinks by 1mu at most, so tight lines keep their operators
\medmuskip=4mu plus 2mu minus 1mu
\thickmuskip=5mu plus 3mu minus 1mu
% headings on the 1.2 scale, semibold
\usepackage{titlesec}
\titleformat*{\section}{\fontsize{15.84}{19}\fontseries{sb}\selectfont}
\titleformat*{\subsection}{\fontsize{13.2}{16}\fontseries{sb}\selectfont}
% no heading below subsection: \subsubsection runs in, like \paragraph
\titleformat{\subsubsection}[runin]{\normalfont\bfseries}{\thesubsubsection}{0.5em}{}[.]
% a part (the complete edition) opens as a title page does: numeral, name,
% subtitle; its title reads "Name: subtitle"
\titleformat{\part}[display]{\raggedright}%
  {\fontsize{47.3}{50}\selectfont\thepart}{6mm}{\arb@parttitle}
\titlespacing*{\part}{0pt}{0pt}{12mm}
\def\arb@parttitle#1{\arb@partsplit#1\@nil}
\def\arb@partsplit#1: #2\@nil{{\fontsize{22.8}{27}\selectfont#1\par}%
  \vspace{3mm}{\fontsize{13.2}{17}\selectfont#2\par}}
% captions small, flush left, bold label; table captions above, wherever the
% source puts them
\usepackage{float}
\floatstyle{plaintop}
\restylefloat{table}
\usepackage[font=small,labelfont=bf,labelsep=period,
  justification=raggedright,singlelinecheck=false]{caption}
\captionsetup[table]{position=top}
% tables are small, set in \@floatboxreset, after the float's \normalsize
\def\arb@table{table}
\g@addto@macro\@floatboxreset{\ifx\@captype\arb@table\small\fi}
% a table's paragraph columns are set ragged, as narrow measures want; \\ still ends
% the row
\g@addto@macro\@arrayparboxrestore{\raggedright\let\\\tabularnewline}
% links: a conifer underline on screen, nothing in print
\hypersetup{pdfborder={0 0 0.6},pdfborderstyle={/S/U/W 0.6},
  allbordercolors={0.110 0.353 0.227},pdflinkmargin=0pt}
% contents: no link underlines; part and section entries in the heading
% weight, ragged right
\patchcmd{\l@part}{\bfseries}{\fontseries{sb}\selectfont}{}{\PatchFailed}
\patchcmd{\l@section}{\bfseries}{\fontseries{sb}\selectfont}{}{\PatchFailed}
\patchcmd{\l@part}{\rightskip\@pnumwidth}{\rightskip\@pnumwidth\@plus1fil}{}{\PatchFailed}
\patchcmd{\l@section}{\rightskip\@pnumwidth}{\rightskip\@pnumwidth\@plus1fil}{}{\PatchFailed}
\patchcmd{\l@part}{\hfil}{\hfill}{}{\PatchFailed}
\patchcmd{\l@section}{\hfil}{\hfill}{}{\PatchFailed}
\let\arb@tableofcontents\tableofcontents
\renewcommand{\tableofcontents}{{\hypersetup{pdfborderstyle={/S/U/W 0}}%
  \arb@tableofcontents}}
% abstract: a flush-left head like the contents', body-size text
\AtBeginDocument{\renewenvironment{abstract}{\section*{\abstractname}}{\par}}
% identity from the title line: name, subtitle, accession numeral
\providecommand{\accession}{}
\def\arb@name#1#2\@nil{\def\arbvolname{#2}}
\def\arb@parsetitle\textbf#1\\[#2]#3#4\@nil{\arb@name#1\@nil\def\arb@sub{#4}}
\providecommand{\arbvolname}{}
\renewcommand{\maketitle}{\expandafter\arb@parsetitle\@title\@nil
  \hypersetup{pdftitle={\arbvolname: \arb@sub},pdfauthor={Marek}}%
  \thispagestyle{empty}\begingroup\raggedright
  {\fontsize{13.2}{16}\selectfont\scshape The Zcash Arboretum\par}
  \vspace*{50mm}
  {\fontsize{47.3}{50}\selectfont\accession\par}
  \vspace{6mm}
  {\fontsize{22.8}{27}\selectfont\arbvolname\par}
  \vspace{3mm}
  {\fontsize{13.2}{17}\selectfont\arb@sub\par}
  \vspace{18mm}
  {\fontsize{13.2}{17}\selectfont Marek\par}
  {\small\@author\par}
  \endgroup\clearpage}
\usepackage{fancyhdr}
\pagestyle{fancy}
\fancyhf{}
\fancyhead[L]{\small\scshape\ifdefempty{\accession}{}{\accession\quad}\arbvolname}
\fancyhead[R]{\small\itshape\nouppercase{\leftmark}}
\fancyfoot[C]{\small\thepage}
\renewcommand{\headrulewidth}{0pt}
\makeatother
\renewcommand{\sectionmark}[1]{\markboth{\thesection\ \ #1}{}}

\newcommand{\F}{\mathbb{F}}
\newcommand{\Z}{\mathbb{Z}}
\newcommand{\N}{\mathbb{N}}
\newcommand{\G}{\mathbb{G}}
\newcommand{\E}{\mathbb{E}}
\newcommand{\PP}{\mathbb{P}}
\renewcommand{\Pr}{\mathrm{Pr}}
\newcommand{\Fp}{\F_p}
\newcommand{\Fq}{\F_q}
\newcommand{\ip}[2]{\langle #1,\, #2 \rangle}
\newcommand{\Comm}{\mathsf{Commit}}
\newcommand{\Open}{\mathsf{Open}}
\newcommand{\Setup}{\mathsf{Setup}}
\newcommand{\Verify}{\mathsf{Verify}}
\newcommand{\negl}{\mathsf{negl}}
\newcommand{\poly}{\mathsf{poly}}
\newcommand{\PPT}{\mathsf{PPT}}
\newcommand{\Adv}{\mathcal{A}}
\newcommand{\Prov}{\mathcal{P}}
\newcommand{\Ver}{\mathcal{V}}
\newcommand{\Sim}{\mathcal{S}}
\newcommand{\Ext}{\mathcal{E}}
\newcommand{\Rel}{\mathcal{R}}
\newcommand{\Lang}{\mathcal{L}}
\newcommand{\nfsym}{\mathsf{nf}}
\newcommand{\cmsym}{\mathsf{cm}}
\newcommand{\cmx}{\mathsf{cmx}}
\newcommand{\cvsym}{\mathsf{cv}}
\newcommand{\rk}{\mathsf{rk}}
\newcommand{\Pallas}{\ensuremath{\mathsf{Pallas}}}
\newcommand{\Vesta}{\ensuremath{\mathsf{Vesta}}}
\newcommand{\Ep}{\ensuremath{\mathcal{E}}}
\newcommand{\NoteCommit}{\ensuremath{\mathsf{NoteCommit}}}
\newcommand{\Sinsemilla}{\ensuremath{\mathsf{Sinsemilla}}}
\newcommand{\ShortCommit}{\ensuremath{\mathsf{ShortCommit}}}
\newcommand{\CommitIvk}{\ensuremath{\mathsf{Commit}^{\mathsf{ivk}}}}
\newcommand{\Extract}{\ensuremath{\mathsf{Extract}}}
\newcommand{\Acc}{\ensuremath{\mathsf{Acc}}}
\newcommand{\GroupHash}{\ensuremath{\mathsf{GroupHash}}}
\newcommand{\ask}{\ensuremath{\mathsf{ask}}}
\newcommand{\aksym}{\ensuremath{\mathsf{ak}}}
\newcommand{\nksym}{\ensuremath{\mathsf{nk}}}
\newcommand{\ivk}{\ensuremath{\mathsf{ivk}}}
\newcommand{\fvk}{\ensuremath{\mathsf{fvk}}}
\newcommand{\ovk}{\ensuremath{\mathsf{ovk}}}
\newcommand{\dk}{\ensuremath{\mathsf{dk}}}
\newcommand{\pkd}{\ensuremath{\mathsf{pk_d}}}
\newcommand{\gd}{\ensuremath{\mathsf{g_d}}}
\newcommand{\rcv}{\ensuremath{\mathsf{rcv}}}
\newcommand{\rcm}{\ensuremath{\mathsf{rcm}}}
\newcommand{\rivk}{\ensuremath{\mathsf{r_{ivk}}}}
\newcommand{\rsk}{\ensuremath{\mathsf{rsk}}}

\renewenvironment{abstract}{%
  \small\begin{center}{\bfseries\abstractname}\end{center}\medskip\noindent\ignorespaces
}{\par\medskip}

% Crosslink notation (one meaning per symbol; see architecture.md, Notation)
% chain kinds, genesis blocks and prefix orders; italic x ranges over the two
% kinds (no star placeholder)
\newcommand{\bc}{\mathrm{bc}}
\newcommand{\bft}{\mathrm{bft}}
\newcommand{\Origin}{\mathcal{O}}
\newcommand{\Obc}{\Origin_{\bc}}
\newcommand{\Obft}{\Origin_{\bft}}
\newcommand{\preceqbc}{\preceq_{\bc}}
\newcommand{\preceqbft}{\preceq_{\bft}}
\newcommand{\precbc}{\prec_{\bc}}
\newcommand{\precbft}{\prec_{\bft}}
\newcommand{\agrees}{\mathrel{\preceq\mkern-9mu\succeq}}
% C truncated by k: C\trunc{k}, or C\trunc[\bc]{k}; an ordinary atom, so a
% following relation keeps its space
\NewDocumentCommand{\trunc}{o m}{\mathord{\lceil}\IfValueT{#1}{_{#1}}^{#2}}
\newcommand{\lca}{\mathsf{lca}}
\newcommand{\score}{\mathsf{score}}
\newcommand{\work}{\mathsf{work}}
\newcommand{\ch}{\mathsf{ch}}
% the sub-protocols and the paper's protocols
\newcommand{\Piorigbc}{\Pi_{\mathrm{origbc}}}
\newcommand{\Piorigbft}{\Pi_{\mathrm{origbft}}}
\newcommand{\Pibc}{\Pi_{\bc}}
\newcommand{\Pibft}{\Pi_{\bft}}
\newcommand{\Pilc}{\Pi_{\mathrm{lc}}}
\newcommand{\Pisac}{\Pi_{\mathrm{sac}}}
% the paper's ledgers; the option is the node: \LOGfin[i]^t
\newcommand{\LOG}{\mathsf{LOG}}
\NewDocumentCommand{\LOGfin}{o}{\LOG_{\mathrm{fin}\IfValueT{#1}{,#1}}}
\NewDocumentCommand{\LOGda}{o}{\LOG_{\mathrm{da}\IfValueT{#1}{,#1}}}
\NewDocumentCommand{\LOGlc}{o}{\LOG_{\mathrm{lc}\IfValueT{#1}{,#1}}}
\NewDocumentCommand{\LOGba}{o}{\LOG_{\mathrm{ba}\IfValueT{#1}{,#1}}}
\newcommand{\sanitize}{\mathsf{sanitize}}
\newcommand{\flatten}{\mathsf{flatten}}
\newcommand{\Ch}{\mathsf{Ch}}
\newcommand{\GST}{\mathsf{GST}}
\newcommand{\GAT}{\mathsf{GAT}}
\newcommand{\Tconfirm}{T_{\mathrm{confirm}}}
\newcommand{\Env}{\mathcal{Z}}
% the BFT interface
\newcommand{\proofsym}{\mathsf{proof}}
\newcommand{\bftlastfinal}{\mathsf{bft\text{-}last\text{-}final}}
\newcommand{\origbftlastfinal}{\mathsf{origbft\text{-}last\text{-}final}}
% the construction
\newcommand{\isstalledblock}{\mathsf{is\text{-}stalled\text{-}block}}
\newcommand{\iscoinbaseonlyblock}{\mathsf{is\text{-}coinbase\text{-}only\text{-}block}}
\newcommand{\contextbft}{\mathsf{context\_bft}}
\newcommand{\headersbc}{\mathsf{headers\_bc}}
\newcommand{\snapshot}{\mathsf{snapshot}}
\newcommand{\LF}{\mathsf{LF}}
\newcommand{\candidate}{\mathsf{candidate}}
\newcommand{\finalitydepth}{\mathsf{finality\text{-}depth}}
% the ledgers: \fin_i^t, (\ba_\mu)_i^t
\newcommand{\fin}{\mathsf{fin}}
\newcommand{\ba}{\mathsf{ba}}
\newcommand{\updatefin}{\mathsf{update\text{-}fin}}
\newcommand{\Bckpt}{B_{\mathrm{ckpt}}}
% lower volumes and ZIPs: the Wallet Guide's confirmation policy, ZIP 237
\newcommand{\ctr}{c_{\mathrm{tr}}}
\newcommand{\cuntr}{c_{\mathrm{untr}}}
\newcommand{\removed}{\mathsf{removed}}


\title{\textbf{\Huge Crosslink in Brief}\\[6pt]\large The core ideas of the Crosslink~2 construction}
\author{m@rek.onl}
\date{}
\begin{document}
\maketitle
\thispagestyle{empty}

\begin{abstract}
This short version of the \emph{Crosslink Guide} states the core of
Crosslink~2, the construction of the design book \emph{Zcash Trailing
Finality Layer} of N.~Wilcox and D.-E.~Hopwood, which adds a
Byzantine-fault-tolerant finality protocol to Zcash's proof-of-work best
chain. It states the problem the construction addresses, its cross-links and
validity rules, the two chains each node derives from them, what is proved
of their safety and what is only argued of liveness, Stalled Mode, the
departures of Shielded Labs' design, the finaliser set, and what ``final''
means to a client. It gives definitions and statements without proofs; each
is cited to the section of the \emph{Crosslink Guide} that treats it and
gives the complete rules and sources. The design runs only on Shielded
Labs' test network. The document assumes the
\emph{Consensus Guide} for work, the best chain and confirmations. It is
non-normative: the Zcash Protocol Specification and the ZIPs are
authoritative, and neither specifies Crosslink.
\end{abstract}
\clearpage \tableofcontents \clearpage

\section{Scope and status}\label{sec:scope}

Crosslink~2 is the construction of N.~Wilcox and D.-E.~Hopwood, \emph{Zcash
Trailing Finality Layer}, called \emph{the book} below, whose Crosslink~2
chapters are unchanged since 13 December 2024; the book calls itself ``an
early and incomplete protocol design proposal''. It is a hybrid of Zcash's
proof-of-work best chain and a \emph{Byzantine-fault-tolerant} (BFT)
protocol: a protocol whose participants decide by votes and which remains
safe while fewer than a fixed fraction of them deviate arbitrarily from it.
The \emph{Crosslink Guide} gives the construction in full, with every rule,
proof and source; this document states its core and cites, for each item,
the section of the \emph{Crosslink Guide} that treats it.

As of 26 October 2026, no ZIP specifies Crosslink. Draft Crosslink texts
exist in an open, unmerged draft pull request to the ZIP repository, each
marked ``ZIP: Unassigned''; they are not ZIPs, they weaken or drop several
rules of the book, and they are not used here. The design runs only on
Shielded Labs' test network, relaunched on 1 October 2026: there is no
deployment on Mainnet or Testnet, no activation decision and no production
value of any parameter. The network upgrade NU7 (ZIP 259, Proposed) does not
include Crosslink. Its 25-second block target spacing (ZIP 218, Proposed)
makes every Crosslink parameter counted in blocks cover a third of its
earlier duration. As the \emph{Crosslink Guide} records as of 16 October
2026 (\S``Scope and status''), no security audit firm has reviewed the
design, and no paper on arXiv or in the IACR ePrint archive analyses it.

Classes and status words are those of the \emph{Crosslink Guide},
\S``Epistemic classes''. A component, that is, a data structure, rule,
parameter or check, is \emph{specified} when normative text of the protocol
specification or of a ZIP fixes it, \emph{designed but unspecified} when a
published design with an analysis exists but no normative text fixes it, and
an \emph{open problem} when no design exists. No Crosslink rule is
specified: every component of Crosslink is designed but unspecified or an
open problem. A result carries one status word: ``proved ($S$)'' when the
source $S$ gives a complete proof, $S$ being the ebb-and-flow paper of Neu,
Tas and Tse (\S\ref{sec:problem}), the TFL book or the \emph{Crosslink
Guide}; ``argued ($S$)'' when $S$ gives a prose argument only;
``assumption''; or ``open problem''. Every Crosslink result below concerns
a designed-but-unspecified construction. A proved result is an implication
about a design and makes no component specified.

\section{The availability-finality dilemma}\label{sec:problem}

Under proof of work, a block gains confirmations as blocks are mined above
it, and the probability that it is reversed falls with their number but
never reaches zero (\emph{Consensus Guide}, \S``Confirmation requirements'',
Proposition ``No finite count is final''). No count of confirmations makes a
block final. A BFT protocol decides by a quorum of votes, and its decisions
are not reversed while fewer than a fixed fraction of its participants, a
third in the protocols considered here, deviate from it; but when too few
participants are online, no quorum forms and the protocol stops.

A \emph{ledger} is a sequence of transactions that a node outputs. A ledger
protocol is \emph{safe} when the ledgers of honest nodes never conflict, one
always being a prefix of the other, and \emph{live} when every transaction
given to honest nodes eventually enters the ledger of every honest node.
Under \emph{dynamic participation}, honest nodes stop and resume taking part
without notice; a \emph{network partition} is an interval during which
messages between two sets of honest nodes are not delivered. The
\emph{availability-finality dilemma} is the statement that no ledger
protocol is both live under dynamic participation and safe under network
partition: a node cannot tell sleeping peers from partitioned ones, so it
either keeps extending its ledger and loses safety under partition, or waits
for a quorum and loses liveness when many nodes sleep. The dilemma is used
here as a statement, not as a theorem (\emph{Crosslink Guide}, \S``The
availability-finality dilemma'').

J.~Neu, E.~N.~Tas and D.~Tse, ``Ebb-and-Flow Protocols: A Resolution of the
Availability-Finality Dilemma'', resolve the dilemma by giving each node two
ledgers: a \emph{finalised} ledger, safe under network partition, and a
\emph{dynamically available} ledger, live under dynamic participation, of
which the finalised ledger is a prefix. Each client chooses which of the two
to trust. Their construction, \emph{snap-and-chat}, runs a longest-chain
protocol and a BFT protocol together. A BFT proposal references a block of
the longest chain, whose chain is the proposal's \emph{snapshot}; an honest
node votes for a proposal only if it sees its snapshot as confirmed in its
own longest chain; and the finalised ledger is the transactions of the
finalised snapshots in order, with duplicates and transactions invalid at
their position removed, a step called sanitisation. The paper proves its
security theorem for a permissioned instance: proved (ebb-and-flow paper).
Its statement for proof of work is supported by one sentence: argued
(ebb-and-flow paper). The paper proves no theorem about Crosslink or about
proof of work (\emph{Crosslink Guide}, \S``The snap-and-chat construction''
and \S``The proof-of-work setting'').

\begin{remark}[From snap-and-chat to Crosslink~2]\label{rem:deficiencies}
Four of the respects in which the book departs from snap-and-chat determine
the construction of Crosslink~2 (\emph{Crosslink Guide}, \S``From
snap-and-chat to Crosslink 2'').
\begin{enumerate}[label=(\arabic*)]
\item \emph{Target.} The finalised ledger is to remain safe even when the
  BFT protocol is entirely subverted, provided the adversary holds less
  than half of the hash rate.
\item \emph{Checkable confirmation.} A proposal carries the headers of a
  fixed number of best-chain blocks above its snapshot, so that any party
  can check from their proof of work that the snapshot is confirmed.
\item \emph{No sanitisation.} Sanitisation can filter the concatenated
  transactions of all snapshots, or remove duplicate snapshot blocks first;
  the two readings coincide when a transaction invalid at a position stays
  invalid at every later one, as for a double spend or a missing input. A
  Zcash shielded transaction names an \emph{anchor}, the root of a pool's
  note commitment tree after some earlier block (\emph{Ironwood Guide},
  \S``Anchors''), and the same root can arise from unrelated transactions,
  so a transaction invalid at its position for a missing anchor can become
  valid at a later one. Both ledgers are therefore plain block chains, and
  no transaction is reordered or removed.
\item \emph{Spending finalised outputs.} Validity of a block may depend only
  on the block and its ancestors, so each best-chain block commits to a BFT
  block from which its finalised prefix is computable, and that commitment
  never moves backwards along a chain.
\end{enumerate}

Status: argued (TFL book).
\end{remark}

\section{The two sub-protocols}\label{sec:sub}

\begin{definition}[Crosslink~2]\label{def:crosslink}
Crosslink~2 keeps a best-chain protocol $\Piorigbc$ and a BFT protocol
$\Piorigbft$ as black boxes and transforms them into protocols $\Pibc$ and
$\Pibft$ by structural additions, changed validity rules and changed rules of
honest behaviour, so that the blocks of each commit to blocks of the other.
These mutual commitments are the \emph{cross-links}. Every node takes part in
both protocols.
\end{definition}

The blocks, chains and genesis blocks of the two protocols are bc-blocks,
bc-chains and $\Obc$, and bft-blocks, bft-chains and $\Obft$. A node $i$ is
\emph{honest at time $t$} in an execution if it has followed the protocol up
to $t$ and the adversary has not cut it off from the network, for example by
an eclipse; an eclipsed node is not honest, and no guarantee below covers
it. No synchrony is assumed unless a sub-protocol imports it. The properties
below are properties of single executions, and the results are implications
between them: when every execution with property $\mathcal{P}$ has property
$\mathcal{Q}$, the probability that $\mathcal{Q}$ fails is at most that of
$\mathcal{P}$ failing. No probabilistic security theorem in the sense of the
ebb-and-flow paper exists for Crosslink~2 (\emph{Crosslink Guide}, \S``The
hybrid construction'', \S``Executions, chains and prefix order'', \S``The
best-chain protocol'' and \S``From snap-and-chat to Crosslink 2'').

A block is identified with the chain of which it is the tip, and the
\emph{height} $\mathsf{height}(B)$ of a block $B$ is its number of
predecessors, so that a genesis block has height~$0$. For a chain $C$ and
$k \in \N$, the chain $C\trunc{k}$, with subscript $\bc$ or $\bft$ where the
kind matters, is the prefix of $C$ whose tip has height
$\max\{0, \mathsf{height}(C) - k\}$: the chain $C$ without its last $k$
blocks. A chain $B$ is a prefix of $C$, written $B \preceq C$ (with subscript
$\bc$ or $\bft$), when $B$ is an initial segment of $C$, equality included;
$B$ and $C$ \emph{agree}, written $B \agrees C$, when one is a prefix of the
other, and \emph{conflict} otherwise. Agreement is not transitive: the
genesis block agrees with two chains that conflict.

\subsection{The best-chain protocol}\label{sec:sub-bc}

The view of a node in $\Piorigbc$ is a set of bc-chains rooted at $\Obc$;
validity of a block depends only on the block and its ancestors; and the
best chain $\ch_i^t$ of node $i$ at time $t$ is a valid chain of highest
score in its view. For Zcash the score is accumulated work and the best
chain is the valid chain of most work (\emph{Consensus Guide}, \S``Work, and
the best chain''). The proof of work of a block is checkable from its header
given a valid parent chain.

\begin{definition}[Prefix Consistency and Prefix Agreement]\label{def:pa}
Let $k \in \N$. An execution of a best-chain protocol has \emph{Prefix
Consistency at depth $k$} iff $\ch_i^t\trunc[\bc]{k} \preceqbc \ch_j^u$ for
all times $t \le u$, every node $i$ honest at $t$ and every node $j$ honest
at $u$, $i = j$ permitted. It has \emph{Prefix Agreement at depth $k$} iff
$\ch_i^t\trunc[\bc]{k} \agrees \ch_j^u\trunc[\bc]{k}$ for all times $t$ and
$u$ and the same nodes.
\end{definition}

Prefix Consistency at depth $k$ implies Prefix Agreement at depth $k$.
Assuming Prefix Consistency at depth $k$ for proof of work excludes network
partitions longer than about $k$ block times: once the best chain on one side
of a partition has grown by more than $k$ blocks, its truncation contains
blocks that the other side cannot know (\emph{Crosslink Guide}, \S``The
best-chain protocol'').

\begin{remark}[The prefix properties for Zcash]\label{rem:pow-prefix}
For Zcash proof of work, Prefix Consistency and Prefix Agreement at any
depth are assumptions, not results. No source proves either property for
Zcash, and no source computes an adequate depth under the 25-second spacing.
Every result below whose premise is a prefix property is therefore
conditional for Zcash.

Status: open problem.
\end{remark}

\subsection{The BFT protocol}\label{sec:sub-bft}

An execution of $\Piorigbft$ proceeds in BFT epochs. In each, a proposer may
make a \emph{proposal}, which names its parent bft-block and its BFT epoch
and is signed by its proposer; for a proposal $P$, as for a bft-block,
$P\trunc[\bft]{1}$ denotes its parent. Each BFT epoch has $n$ \emph{voting
units},
distributed among the \emph{validators} of the BFT protocol; the holders of
a BFT epoch's units form its \emph{finaliser set}, each being a
\emph{finaliser}. A proposal is \emph{notarised} by ballots cast for it by
at least $q = \lceil 2n/3 \rceil$ units of its BFT epoch, a protocol being
free to require more, and a bft-block is a notarised proposal together with
its notarisation proof. By the \emph{Parent rule}, a bft-block other than
$\Obft$ is valid only if its parent is valid.

\begin{definition}[Finality in context]\label{def:blf}
The protocol supplies an efficiently computable function $\bftlastfinal$
that maps each valid bft-block $C$ to a prefix $\bftlastfinal(C)
\preceqbft C$, the \emph{last final block in the context of $C$}, with
$\bftlastfinal(\Obft) = \Obft$ and $\bftlastfinal(C) \preceqbft
\bftlastfinal(D)$ whenever $C \preceqbft D$ are valid. A bft-block $B$ is
\emph{final in the context of $C$} iff $B \preceqbft \bftlastfinal(C)$. No
bft-block is final other than in a context.
\end{definition}

\begin{assumption}[One-third bound]\label{ass:one-third}
In the execution considered, for every BFT epoch, strictly fewer than one
third of its voting units are ever cast non-honestly. A unit is cast
non-honestly when a party other than its holder casts it, for example after
a key compromise, when it is cast for two distinct proposals of the BFT
epoch, or when its holder casts it against the rules of honest voting.

Status: assumption.
\end{assumption}

The bound covers all time: a key compromised long after its BFT epoch can
still cast that BFT epoch's units.

\begin{definition}[Final Agreement]\label{def:fa}
An execution of a BFT protocol has \emph{Final Agreement} iff
$\bftlastfinal(C) \agrees \bftlastfinal(C')$ for all valid bft-blocks $C$
and $C'$ in honest view at any times.
\end{definition}

If Assumption~\ref{ass:one-third} fails at some time, Final Agreement may
fail from that time on. Unlike the prefix properties of proof of work, Final
Agreement can hold in an asynchronous network (\emph{Crosslink Guide},
\S``The BFT protocol'').

\section{The construction}\label{sec:cons}

\begin{definition}[Parameters]\label{def:params}
Crosslink~2 has a \emph{confirmation depth} $\sigma \in \N^{+}$, a
\emph{finalisation gap bound} $L \in \N^{+}$, and a predicate
$\isstalledblock$ on bc-blocks. No value of $\sigma$ or $L$ is fixed: the
values are an open problem.
\end{definition}

The book requires $L$ to be ``significantly greater than $\sigma$'' and
``\emph{at least} $2\sigma$'' in practice, without derivation; the
\emph{Crosslink Guide} proves that $L \ge \sigma + 1$ is necessary for
Stalled Mode to end while the BFT protocol is live (\S\ref{sec:live}).

\begin{definition}[Cross-link fields]\label{def:fields}
Each bc-header carries a field $\contextbft$, a commitment to a bft-block,
the \emph{context} of the block. Each proposal and each bft-block $B$ other
than $\Obft$ carries a field $B.\headersbc$, a sequence of $\sigma$
bc-headers, deepest first. The \emph{$\sigma$-block tail} of a bc-chain $T$
of at least $\sigma + 1$ blocks is the sequence of the headers of
$T\trunc{\sigma-1}, \ldots, T\trunc{1}, T$.
\end{definition}

\begin{definition}[Derived quantities]\label{def:derived}
For a proposal or bft-block $B$, the \emph{snapshot} $\snapshot(B)$ is the
parent of the deepest header that $B$ carries, and $\snapshot(\Obft) = \Obc$;
when $B$ carries the $\sigma$-block tail of $T$,
$\snapshot(B) = T\trunc{\sigma}$.
For a bc-block $H$,
\begin{align*}
  \LF(H) &:= \bftlastfinal(H.\contextbft),\\
  \finalitydepth(H) &:= \mathsf{height}(H) -
    \mathsf{height}\bigl(\snapshot(\LF(H))\bigr),\\
  \candidate(H) &:= \lca\bigl(\snapshot(\LF(H)),\, H\trunc{\sigma}\bigr),
\end{align*}
where $\lca$ is the last common ancestor. The block $\LF(H)$ is the
\emph{last final context} of $H$, and $\snapshot(\LF(H))$ is its \emph{final
snapshot}.
\end{definition}

A bft-block thus points down into the best chain through its carried
headers, and a bc-block points up into the BFT chain through its context
(Figure~\ref{fig:crosslinks}; \emph{Crosslink Guide}, \S``Parameters'' and
\S``Structural additions''); each node computes its finalised chain from
these quantities (\S\ref{sec:ledgers}).

\begin{figure}[ht]
\centering
\begin{tikzpicture}[x=12mm, y=10mm,
  blk/.style={draw, rectangle, minimum width=5.5mm, minimum height=5.5mm,
    inner sep=0pt, fill=white},
  fin/.style={blk, fill=tierin!25}, hdr/.style={blk, fill=tierfull!30},
  dts/.style={inner sep=1pt}, lnk/.style={-{Stealth[length=2mm]}, thick},
  snp/.style={-{Stealth[length=2mm]}, dashed},
  ctx/.style={-{Stealth[length=2mm]}, thick, tierrecv},
  lbl/.style={font=\footnotesize},
]
% bc-chain
\node[blk] (g) at (0,0) {};
\node[dts] (d1) at (1,0) {$\cdots$};
\node[fin] (s1) at (2,0) {};
\node[dts] (d2) at (3,0) {$\cdots$};
\node[blk] (s2) at (4,0) {};
\node[hdr] (h0) at (5,0) {};
\node[dts] (d3) at (6,0) {$\cdots$};
\node[hdr] (hs) at (7,0) {};
\node[dts] (d4) at (8,0) {$\cdots$};
\node[blk] (H) at (9,0) {};
\draw[lnk] (d1) -- (g);
\draw[lnk] (s1) -- (d1);
\draw[lnk] (d2) -- (s1);
\draw[lnk] (s2) -- (d2);
\draw[lnk] (h0) -- (s2);
\draw[lnk] (d3) -- (h0);
\draw[lnk] (hs) -- (d3);
\draw[lnk] (d4) -- (hs);
\draw[lnk] (H) -- (d4);
\node[lbl, below=1mm of g] {$\Obc$};
\node[lbl, below=1mm of s1, align=center]
  {$\snapshot(B_1)$\\$=\snapshot(\LF(H))$};
\node[lbl, below=1mm of s2] {$\snapshot(B_2)$};
\node[lbl, below=1mm of h0] {$[0]$};
\node[lbl, below=1mm of hs] {$[\sigma - 1]$};
\node[lbl, below=1mm of H] {$H$};
% the sigma carried headers of B2
\draw (4.7,0.45) -- (4.7,0.6) -- (7.3,0.6) -- (7.3,0.45);
\node[lbl, above] (hl) at (6,0.6) {$B_2.\headersbc$};
% bft-chain
\node[blk] (o) at (1,2.4) {};
\node[fin] (b1) at (3.5,2.4) {};
\node[blk] (b2) at (6,2.4) {};
\draw[lnk] (b1) -- (o);
\draw[lnk] (b2) -- (b1);
\node[lbl, above=1mm of o] {$\Obft$};
\node[lbl, above=1mm of b1] {$B_1 = \LF(H)$};
\node[lbl, above=1mm of b2] {$B_2$};
\draw[snp] (o) -- (g);
\draw[snp] (b1) -- (s1);
\draw[snp] (b2) -- (s2);
\draw[dotted] (b2) -- (hl);
\draw[ctx] (H) to[out=90, in=0]
  node[lbl, pos=0.4, right=1mm] {$H.\contextbft$} (b2);
\end{tikzpicture}
\caption{The cross-links. Solid black arrows point from child to parent on
the bc-chain (below) and the bft-chain $\Obft$, $B_1$, $B_2$ (above); dashed
arrows point from a bft-block to its snapshot. The bracket marks the
$\sigma$ headers $B_2.\headersbc[0]$ to $B_2.\headersbc[\sigma - 1]$ (orange)
that $B_2$ carries, and $\snapshot(B_2)$ is the parent of the deepest of
them. The green arrow is $H.\contextbft = B_2$. In the case drawn, the last
block final in the context of $B_2$ is $B_1$ (blue), so $\LF(H) = B_1$ and
the final snapshot of $H$ is $\snapshot(B_1)$ (blue). The Linearity rule
appears as $\snapshot(B_1)$ lying at or below $\snapshot(B_2)$ on one
bc-chain.}
\label{fig:crosslinks}
\end{figure}

\subsection{The rules}\label{sec:cons-rules}

\begin{definition}[Validity in $\Pibft$]\label{def:bft-validity}
A proposal or bft-block $B$ other than $\Obft$ is valid in $\Pibft$ iff it
is valid in $\Piorigbft$ and
\begin{enumerate}[label=(\roman*)]
\item \emph{Linearity rule}: $\snapshot(B\trunc[\bft]{1}) \preceqbc
  \snapshot(B)$, so that snapshots never move backwards along a bft-chain;
\item \emph{Tail Confirmation rule}: $B.\headersbc$ is the $\sigma$-block
  tail of a bc-valid chain.
\end{enumerate}
The function $\bftlastfinal$ of $\Pibft$ is that of $\Piorigbft$.
\end{definition}

\begin{definition}[Honest voting]\label{def:honest-voting}
A validator $v$ at time $t$ first adds the bc-blocks of the carried headers
of a proposal $P$ to its view, checking their validity, and then votes for
$P$ only if $P$ is valid in $\Pibft$ and the \emph{Confirmed best-chain
criterion} holds:
\[
  \snapshot(P) \preceqbc \ch_v^t\trunc{\sigma},
\]
that is, the snapshot is $\sigma$-confirmed in the validator's own best
chain.
\end{definition}

\begin{definition}[Validity in $\Pibc$]\label{def:bc-validity}
The genesis block has $\Obc.\contextbft = \Obft$. A bc-block $H \neq \Obc$ is
valid in $\Pibc$ iff it is valid in $\Piorigbc$ and
\begin{enumerate}[label=(\roman*)]
\item \emph{Valid context rule}: $H.\contextbft$ is a valid bft-block;
\item \emph{Extension rule}: $\LF(H\trunc[\bc]{1}) \preceqbft \LF(H)$, so
  that the last final context never moves backwards along a bc-chain;
\item \emph{Last Final Snapshot rule}: $\snapshot(\LF(H)) \preceqbc H$, so
  that the final snapshot lies on the block's own chain;
\item \emph{Finality depth rule}: $\finalitydepth(H) \le L$ or
  $\isstalledblock(H)$.
\end{enumerate}
The rules on transactions are those of $\Piorigbc$, unchanged.
\end{definition}

\begin{definition}[Unchanged fork choice]\label{def:fork-choice}
An honest node uses no information from $\Pibft$ other than the validity
rules of Definition~\ref{def:bc-validity} to choose the bc-chain it follows:
$\ch_i^t$ is a bc-valid chain of highest score in its view, as in
$\Piorigbc$.
\end{definition}

The BFT protocol steers the best chain only through validity: a chain whose
blocks cite contexts that move finality backwards, or place the final
snapshot off the chain, is invalid, and the most-work rule chooses among the
rest. The rules answer items~2 to~4 of Remark~\ref{rem:deficiencies}: the
Tail Confirmation rule makes confirmation checkable by any party (item~2),
while the Confirmed best-chain criterion keeps snap-and-chat's voting
condition; the Linearity rule makes the snapshots of a valid
bft-chain extend one another, so that the finalised history is one bc-chain
and needs no sanitisation (item~3); and the Extension and Last Final
Snapshot rules make the finalised prefix computable from a block's own chain
(item~4). The Linearity rule also carries the safety argument of
\S\ref{sec:safety-fa}, and the Finality depth rule bounds availability
(\S\ref{sec:cons-stalled}). All of these
definitions are designed but unspecified (\emph{Crosslink Guide},
\S``Validity rules of the BFT protocol'' to \S``Honest block production and
fork choice'').

\subsection{Stalled Mode}\label{sec:cons-stalled}

\begin{definition}[Stalled Mode]\label{def:stalled}
A bc-block $H$ is in \emph{Stalled Mode} iff $\finalitydepth(H) > L$. By
the Finality depth rule, every valid bc-block in Stalled Mode is
\emph{stalled}: it satisfies $\isstalledblock$.
\end{definition}

For Zcash the book proposes that a stalled block be \emph{coinbase-only}: it
has exactly one transaction, which only distributes newly issued funds and
has no inputs; the book never fixes the predicate. Stalled Mode is an
objective property of each block, not a network-wide switch. In Stalled Mode
the best chain can grow only by stalled blocks, which under the proposed
policy carry no ordinary spends, and that it keeps growing is argued (TFL
book) (\S\ref{sec:live}); a block that cites a context whose
final snapshot lies on its chain at most $L$ blocks below it is not in
Stalled Mode, whatever its parent (\emph{Crosslink Guide}, \S``Stalled
Mode'').

\section{The two chains of a node}\label{sec:ledgers}

From its best chain, each node derives two chains, which play the roles of
the finalised and the available ledger of the ebb-and-flow model
(\emph{Crosslink Guide}, \S``The locally finalised chain'' and \S``The
locally bounded-available chain'').

\begin{definition}[Locally finalised chain]\label{def:fin}
Each node $i$ keeps a bc-chain $\fin_i^t$, starting at $\fin_i^0 = \Obc$. At
each time $t$ at which its best chain changes, with $s$ the time of the
previous change, or $0$ if there is none,
\[
  \fin_i^t :=
  \begin{cases}
    \candidate(\ch_i^t), & \text{if } \fin_i^s \preceqbc
      \candidate(\ch_i^t),\\
    \fin_i^s, & \text{otherwise};
  \end{cases}
\]
if $\candidate(\ch_i^t)$ conflicts with $\fin_i^s$, the node records a
\emph{finalisation safety hazard}. Between changes of its best chain,
$\fin_i$ is constant.
\end{definition}

\begin{lemma}[Local finalization linearity]\label{lem:fin-linear}
A node that follows the update rule of Definition~\ref{def:fin} has
$\fin_i^t \preceqbc \fin_i^u$ for all $t \le u$: its finalised chain never
rolls back.

Status: proved (Crosslink Guide).
\end{lemma}

The chain $\fin_i^t$ is state of the node, not a field of any block, and no
consensus rule reads it. A BFT decision alone changes nothing at a node. For
a node $i$ honest at $t$, a bc-block $X \neq \Obc$ lies in $\fin_i^t$ exactly
when four conditions hold: $X$ lies in the snapshot of a proposal; the
bft-block formed from that proposal is final in the context of some valid
bft-block; some valid bc-block $H$ with $X \preceqbc H\trunc{\sigma}$ has a
last final context that extends that bft-block; and $H$ has become the
node's best tip, by time $t$, at an update that set $\fin_i$ to
$\candidate(H)$ (\emph{Crosslink Guide}, \S``Finality as a status'',
Proposition ``Path to finality''). When $\LF(H) \neq \Obft$ and the headers
carried by $\LF(H)$ are ancestors of $H$, $\candidate(H)$ is the final
snapshot of $H$, so an update
to it leaves $\fin_i$ trailing the tip by $\finalitydepth(H) \ge \sigma + 1$
blocks; the bound is attained (\emph{Crosslink Guide}, \S``Computation of
the finalised prefix'', Proposition ``Finalisation lag'').

\begin{definition}[Locally bounded-available chain]\label{def:ba}
Each node chooses a depth $\mu$ with $0 < \mu \le \sigma$, by default $\mu =
\sigma$. Its \emph{locally bounded-available chain} at time $t$ is
\[
  (\ba_\mu)_i^t :=
  \begin{cases}
    \ch_i^t\trunc{\mu}, & \text{if } \fin_i^t \preceqbc \ch_i^t\trunc{\mu},\\
    \fin_i^t, & \text{otherwise}.
  \end{cases}
\]
\end{definition}

\begin{proposition}[Ledger prefix property]\label{prop:ledger-prefix}
For every node $i$ honest at time $t$, $\fin_i^t \preceqbc (\ba_\mu)_i^t$.

Status: proved (TFL book).
\end{proposition}

\begin{figure}[ht]
\centering
\begin{tikzpicture}[x=11mm, y=10mm,
  blk/.style={draw, rectangle, minimum width=5.5mm, minimum height=5.5mm,
    inner sep=0pt, fill=white},
  fin/.style={blk, fill=tierin!30}, stl/.style={blk, fill=tierspend!30},
  dts/.style={inner sep=1pt}, lnk/.style={-{Stealth[length=2mm]}, thick},
  lbl/.style={font=\footnotesize}, bar/.style={line width=2.5pt},
]
\node[blk] (g) at (0,0) {};
\node[dts] (d1) at (1,0) {$\cdots$};
\node[fin] (s) at (2,0) {};
\node[dts] (d2) at (3,0) {$\cdots$};
\node[blk] (l) at (4,0) {};
\node[stl] (f) at (5,0) {};
\node[dts] (d3) at (6,0) {$\cdots$};
\node[stl] (ts) at (7,0) {};
\node[dts] (d4) at (8,0) {$\cdots$};
\node[stl] (tm) at (9,0) {};
\node[dts] (d5) at (10,0) {$\cdots$};
\node[stl] (H) at (11,0) {};
\draw[lnk] (d1) -- (g);
\draw[lnk] (s) -- (d1);
\draw[lnk] (d2) -- (s);
\draw[lnk] (l) -- (d2);
\draw[lnk] (f) -- (l);
\draw[lnk] (d3) -- (f);
\draw[lnk] (ts) -- (d3);
\draw[lnk] (d4) -- (ts);
\draw[lnk] (tm) -- (d4);
\draw[lnk] (d5) -- (tm);
\draw[lnk] (H) -- (d5);
\node[lbl, below=1mm of g] {$\Obc$};
\node[lbl, below=1mm of s, align=center]
  {$\snapshot(\LF(H))$\\$= \candidate(H)$};
\node[lbl, below=1mm of ts] {$H\trunc{\sigma}$};
\node[lbl, below=1mm of tm] {$H\trunc{\mu}$};
\node[lbl, below=1mm of H] {$H$};
% brackets above the chain
\draw (1.7,0.45) -- (1.7,0.6) -- (4.3,0.6) -- (4.3,0.45);
\node[lbl, above] at (3,0.6) {$L$};
\draw (6.7,0.45) -- (6.7,0.6) -- (11.3,0.6) -- (11.3,0.45);
\node[lbl, above] at (9,0.6) {$\sigma$};
\node[lbl, text=tierspend] at (5.5,0.75) {stalled};
\draw (1.7,1.25) -- (1.7,1.4) -- (11.3,1.4) -- (11.3,1.25);
\node[lbl, above] at (6.5,1.4)
  {$\finalitydepth(H) \ge \sigma + 1$};
% the two chains of node i
\draw[bar, tierin] (0,-1.6) -- (2,-1.6);
\node[lbl, anchor=west] at (2.1,-1.6) {$\fin_i^t$};
\draw[bar, tierrecv] (0,-2.0) -- (9,-2.0);
\node[lbl, anchor=west] at (9.1,-2.0) {$(\ba_\mu)_i^t$};
\end{tikzpicture}
\caption{The two chains of node $i$ at an update with best tip
$H = \ch_i^t$ in Stalled Mode, when the headers that $\LF(H)$ carries are
ancestors of $H$ and the update sets $\fin_i^t$ to $\candidate(H)$. Arrows
point from child to parent. The final snapshot $\snapshot(\LF(H))$ is the
lower of itself and $H\trunc{\sigma}$, hence $\candidate(H)$; $\fin_i^t$
(blue) ends there, $\finalitydepth(H) \ge \sigma + 1$ blocks below $H$. The
chain $(\ba_\mu)_i^t$ (green) is $H\trunc{\mu}$, and $\fin_i^t$ is a prefix of
it. The blocks more than $L$ above the final snapshot, up to $H$, are stalled
(red). Positions are symbolic and represent no value of $\sigma$, $\mu$ or
$L$.}
\label{fig:two-ledgers}
\end{figure}

\section{Safety}\label{sec:safety}

\begin{definition}[Assured Finality]\label{def:af}
An execution has \emph{Assured Finality} iff $\fin_i^t \agrees \fin_j^u$ for
all times $t$ and $u$, every node $i$ honest at $t$ and every node $j$
honest at $u$, $i = j$ permitted.
\end{definition}

Assured Finality is the book's main safety goal: the finalised chains of
honest nodes may differ in length, but never conflict (\emph{Crosslink
Guide}, \S``Assured Finality''). Two routes lead to it, one through each
sub-protocol.

\subsection{The Prefix Agreement route}\label{sec:safety-pa}

\begin{theorem}[Prefix Agreement implies Assured Finality]\label{thm:pa-af}
In every execution of Crosslink~2 in which $\Pibc$ has Prefix Agreement at
depth $\sigma$, Assured Finality holds.

Status: proved (TFL book).
\end{theorem}

\begin{corollary}[Safety under subversion of the BFT protocol]
\label{cor:bft-subversion}
In every execution in which $\Pibc$ has Prefix Agreement at $\sigma$,
Assured Finality holds whatever the holders of voting units do, including
executions in which every voting unit is cast non-honestly and signatures
are forged.

Status: proved (Crosslink Guide).
\end{corollary}

The proof of Theorem~\ref{thm:pa-af} uses nothing of the BFT protocol: each
finalised chain of a node is a prefix of one of its earlier best chains
truncated by $\sigma$, and Prefix Agreement makes those truncations agree.
This meets the target of item~1 of Remark~\ref{rem:deficiencies},
conditionally on Prefix Agreement of $\Pibc$, whose transfer from
$\Piorigbc$ to $\Pibc$ is open (\emph{Crosslink Guide}, \S``The Prefix
Agreement route'').

\subsection{The Final Agreement route}\label{sec:safety-fa}

\begin{theorem}[Final Agreement implies Assured Finality]\label{thm:fa-af}
In every execution of Crosslink~2 in which $\Pibft$ has Final Agreement,
Assured Finality holds.

Status: proved (Crosslink Guide).
\end{theorem}

The book states this theorem, but its printed proof is a copy of its proof
of Theorem~\ref{thm:pa-af}; the \emph{Crosslink Guide} gives a proof, which
uses neither $\sigma$ nor Prefix Agreement but needs the Linearity rule
(\emph{Crosslink Guide}, \S``The Final Agreement route'').

\begin{proposition}[Necessity of the Linearity rule]\label{prop:fa-needs-lin}
Let the Linearity rule be removed. For every $\sigma$ there is an execution
in which $\Pibft$ has Final Agreement, the best chains of two honest nodes
$i$ and $j$ satisfy the Valid context, Extension and Last Final Snapshot
rules, the voters of every bft-block they cite meet the Confirmed best-chain
criterion, and $\fin_i$ and $\fin_j$ conflict.

Status: proved (Crosslink Guide).
\end{proposition}

\subsection{Consequences}\label{sec:safety-either}

\begin{corollary}[Assured Finality under either premise]\label{cor:af-either}
Assured Finality holds in every execution in which $\Pibc$ has Prefix
Agreement at $\sigma$ or $\Pibft$ has Final Agreement.

Status: proved (Crosslink Guide).
\end{corollary}

The finalised chains of two honest nodes can therefore conflict only in an
execution in which both premises fail. A partition of the network, which
breaks Prefix Agreement, together with an adversary holding more than a
third of the units of a BFT epoch and voting with each side, which breaks
Final Agreement, is such an execution; the book calls this attack
unavoidable, since no best-chain protocol keeps finality under partition
(argued (TFL book); \emph{Crosslink Guide}, \S``Consequences of the two
routes''). For Zcash, both premises are open: Prefix Agreement of proof of
work is an assumption (Remark~\ref{rem:pow-prefix}), and no concrete BFT
protocol, hence no proof of its Final Agreement, is fixed; that Final
Agreement of $\Piorigbft$ carries over to $\Pibft$ is argued (TFL book)
(\emph{Crosslink Guide}, \S``Liveness of the sub-protocols'', Proposition
``Transfer by erasure''). Crosslink is thus
not proved safe for Zcash; it is proved safe under either of two premises,
each of which is open for Zcash.

The bounded-available chain is less safe than the finalised chain.
Agreement of the chains $\ba_\mu$ of honest nodes follows from Prefix
Agreement at depth $\mu$ (proved (Crosslink Guide)), and not from Final
Agreement (\emph{Crosslink Guide}, \S``Safety of the bounded-available
chain'').

\section{Liveness and bounded availability}\label{sec:live}

No liveness statement of Crosslink~2 is proved. The book argues that the
best chain of $\Pibc$ grows whenever that of $\Piorigbc$ does, given that a
coinbase-only stalled block can always be produced; that $\Pibft$ is live
whenever $\Piorigbft$ is; and that Stalled Mode is not entered while there
are enough honest block producers, proposers and validators and $L$ is large
enough: argued (TFL book), each. The transfer of the prefix properties from
$\Piorigbc$ to $\Pibc$, the advance of finalisation and a rule for choosing
$L$ are open problems. Three proved results bound what any liveness argument
can claim (\emph{Crosslink Guide}, \S``Liveness of the sub-protocols'',
\S``Progress of finalisation'' and \S``Stalled Mode''):
\begin{enumerate}[label=(\roman*)]
\item if $L \le \sigma$, Stalled Mode persists even while $\Pibft$ is live,
  so $L \ge \sigma + 1$ is necessary;
\item block producers that never advance $\contextbft$ over $L$ consecutive
  heights force a stalled block whatever $\Pibft$ does;
\item holders of more than one third of the units of a BFT epoch who
  withhold their votes, and whose units no one else casts, prevent every
  notarisation of that BFT epoch.
\end{enumerate}

\emph{Bounded availability} is the book's informal weakening of dynamic
availability: during a finalisation stall, the number of best-chain blocks
above the final snapshot of the best tip may grow only up to a bound,
beyond which ordinary spends stop. The book argues for it that under unbounded
availability the cost of rolling back spent funds grows without limit while
finality stalls, and that halting the chain instead would give an attacker
unbounded time to build a fork of the recent tail: argued (TFL book)
(\emph{Crosslink Guide}, \S``Bounded availability''). The book does not
prove the bound; the \emph{Crosslink Guide} proves it for a node whose
finalised chain lies on its best chain.

\begin{proposition}[Bounded non-stalled run]\label{prop:bounded-run}
Let $H = \ch_i^t$, and let node $i$ have followed the update rule of
Definition~\ref{def:fin} up to $t$. Suppose that the Extension, Last Final
Snapshot and Finality depth rules hold along $H$, that the Parent and
Linearity rules hold along the bft-chain that ends at $\LF(H)$, that $\sigma
\le L$, and that $\fin_i^t \preceqbc H$. Then every block of $H$ at height
greater than $\mathsf{height}(\fin_i^t) + L$ is stalled.

Status: proved (Crosslink Guide).
\end{proposition}

For some best-chain protocol with the properties of \S\ref{sec:sub-bc} the
constraint $\sigma \le L$ cannot be dropped (proved (Crosslink Guide)); the
case in which $\fin_i^t$ is not a prefix of the best chain is open
(\emph{Crosslink Guide}, \S``The bounded non-stalled run''). A \emph{finality
override}, a deliberate reversal of finalised blocks in response to a failure
of the protocol, has no rule, message or state in the construction: it is a
change of the consensus rules, a hard fork outside the protocol, and Assured
Finality binds executions of the protocol only. The book's case for
formalising overrides is argued (TFL book) (\emph{Crosslink Guide},
\S``Finality overrides'').

\section{The Shielded Labs design}\label{sec:sl}

The design that Shielded Labs is building, described in its \emph{Zcash
Crosslink Design Overview} (draft~10, 14 September 2026), keeps the Tail
Confirmation, Linearity, Valid context, Extension and Last Final Snapshot
rules and the update rule of $\fin$. It drops the Finality depth rule,
Stalled Mode, the bound $L$ and the chain $\ba_\mu$. It adds a \emph{sticky
fork choice}, under which a node follows the heaviest valid chain that
contains its finalised chain; an \emph{Inclusion Depth rule}, a lower bound
$\finalitydepth(H) \ge \sigma + 1$, not an upper bound; the BFT protocol
Tenderlink, derived from Tendermint and without a specification, whose
bft-blocks, called \emph{certificates}, carry no signatures: a best-chain
block cites a certificate through a reference that carries its notarisation
signatures; and \emph{slash forks}, hard forks that node operators may adopt
to punish named finalisers (\S\ref{sec:fs}). Every component is designed
but unspecified (\emph{Crosslink Guide}, \S``The Shielded Labs design'').

Under its rules, both routes to Assured Finality are proved (Crosslink
Guide), but neither premise is: Final Agreement of Tenderlink is an open
problem, and Prefix Agreement of the best chain under sticky fork choice is
unproved, while one formal model exhibits its failure on a small instance in
which more than a third of the BFT voting power is subverted. Sticky fork
choice never records a finalisation hazard (proved (Crosslink Guide)); its
safety and liveness in general are open (\emph{Crosslink Guide},
\S``Sticky fork choice''). Without the Finality depth rule,
the run of best-chain blocks above $\fin$ is unbounded: for every $d$ there
is an execution in which the best chain extends $d$ blocks beyond the
finalised chain (proved (Crosslink Guide); \emph{Crosslink Guide},
\S``Theorems under the Shielded Labs design'').

\section{The finaliser set}\label{sec:fs}

The book models only $n$ voting units per BFT epoch and a two-thirds
threshold, and lists the selection of a proof-of-stake sub-protocol and the
rewards of stakers as missing. Every concrete element of the finaliser set
below is Shielded Labs' design and is designed but unspecified, with the open
problems noted (\emph{Crosslink Guide}, \S``The finaliser set'').
\begin{itemize}
\item \emph{Voting power.} The finalisers are the participants of
  Tenderlink. A finaliser's voting power is the stake delegated to it; the
  $N$ finalisers of largest power form the active roster, and a certificate
  needs two thirds of the active roster's stake. Every fraction is a
  fraction of stake, never of the number of finalisers. The value of $N$
  and the selection rule are open.
\item \emph{Roster.} The roster of each BFT height, the height of a
  certificate in the bft-chain, derives from the staking state at the
  snapshot of the certificate decided at the preceding BFT height. That
  snapshot generally lies above a node's finalised chain and may lie off
  its best chain; the roster cannot change while no new certificate is
  decided, except by a slash fork.
\item \emph{Delegation.} A staker locks shielded ZEC into a bond, whose
  amount is a power of ten ZEC, delegated to one finaliser and redelegable
  at any time. Leaving takes two steps, unbonding and withdrawing, in
  different \emph{staking windows}, recurring ranges of best-chain heights to
  which the creation of bonds, unbonding and withdrawal are confined, so
  that misbehaviour can be punished before the funds leave.
\item \emph{Privacy.} Batching in staking windows and amounts quantised at
  creation are meant to keep an observer from linking a bond to the shielded
  notes that funded it. The totals of each finaliser and the amount and
  target of each bond are public; hiding the delegation itself is an open
  problem.
\item \emph{Rewards.} Staking is paid from the existing block subsidy, with
  no new issuance and no change to the distribution of fees. The published
  split gives 40 per cent of the subsidy to miners, 40 per cent to stakers
  and 20 per cent to Zcash's development fund, which is not a Crosslink
  share; of the stakers' part, 90 per cent goes to bonds and 10 per cent to
  finalisers as commission. Whether payment depends on new finality is open.
\item \emph{Slashing.} No rule of the protocol punishes a finaliser. Node
  operators may instead adopt a slash fork, a hard fork that burns the
  bonds delegated to named finalisers and removes those finalisers from the
  roster.
\end{itemize}

\begin{proposition}[Power of a finaliser coalition]\label{prop:coalition}
When ballots cannot be forged and notarisation proofs are sound, holders of
more than one third of a BFT epoch's units can prevent every notarisation of
that BFT epoch by withholding. In every execution, whatever the holders of
voting units do, they cannot make the finalised chains of honest nodes
conflict while $\Pibc$ has Prefix Agreement at $\sigma$.

Status: proved (Crosslink Guide).
\end{proposition}

\section{What ``final'' means to a client}\label{sec:client}

\begin{definition}[Final]\label{def:final}
A bc-block $X$ is \emph{final at node $i$ at time $t$} iff $X \preceqbc
\fin_i^t$. A transaction is final at $i$ at $t$ iff a block that contains
it is final at $i$ at $t$.
\end{definition}

This is the book's assured finality: the protocol cannot revert a final
transaction, while a fork made outside the protocol, a finality override
among them, still can. Finality of a bc-block is relative to one node, and
differs from finality of a bft-block in a context (Definition~\ref{def:blf};
\emph{Crosslink Guide}, \S``Finality as a status'').

\begin{proposition}[Guarantee of finality]\label{prop:final-guarantee}
Let $X$ be final at node $i$ at time $t$, with $i$ honest at $t$. Then $X$ is
final at $i$ at every later time at which $i$ is honest. In every execution
in which $\Pibc$ has Prefix Agreement at $\sigma$ or $\Pibft$ has Final
Agreement, moreover, $X$ and $\fin_j^u$ agree for every time $u$ and every
node $j$ honest at $u$.

Status: proved (Crosslink Guide).
\end{proposition}

At each time, a block is in exactly one of three states at a node: final; not
yet final, when it strictly extends the node's finalised chain; or never
final, when it conflicts with it. The state is decided by ancestry, not by
height, and a block not yet final has no promise of becoming final
(\emph{Crosslink Guide}, \S``Finality as a status'', Proposition ``Finality
status'').

Finality is distinct from confirmation depth. A deep block outside
$\fin_i^t$ is not final, and a block of $\fin_i^t$ is final at any depth;
confirmation counts still measure proof-of-work depth. A block that becomes
final at an update under the hypotheses of the lag bound of
\S\ref{sec:ledgers} has at least $\sigma + 2$ confirmations at that time
(proved (Crosslink Guide); \emph{Crosslink Guide}, \S``Finality and
confirmation depth'').

An application reads one of two views, $\fin_i^t$ or $(\ba_\mu)_i^t$, never
the best chain itself. The choice sets its behaviour during a finalisation
stall: with $\fin$ it stops at once; with $\ba_\mu$ it keeps processing user
transactions until $L$ blocks have passed. Shielded Labs' design has neither
$\ba_\mu$ nor $L$: an application there either stops with $\fin$ or follows
the best chain, which can run unboundedly ahead of it (\emph{Crosslink
Guide}, \S``Ledger views and stalls''). A light client that verifies block
headers only cannot verify finality in any current design: no design places
in the header chain both the notarisation signatures and a commitment to
the finalisers against which they are counted. This is an open problem
(\emph{Crosslink Guide}, \S``Light clients'').

\section{Open problems}\label{sec:open}

A specification of Crosslink~2 for Zcash must fix each of the following
components, and no source fixes any of them (\emph{Crosslink Guide},
\S``Components a specification must fix''): the concrete BFT protocol, the
assignment of its voting units and the format of its notarisation proofs;
the values of $\sigma$ and $L$, with $\sigma + 1 \le L$; the predicate
$\isstalledblock$; the encoding and placement of the cross-link fields; the
mapping from stake to voting units and the selection of the roster; the
staking calendar; the format of bond transactions, the quantised amounts and
a minimum stake; the reward rule and the division of the commission; the
fees of staking transactions; the rotation of finaliser keys and the initial
roster; and activation on the existing chain. Four partial formal models
exist, in Lean, Quint and TLA$^{+}$; none machine-checks Assured Finality or
either theorem of \S\ref{sec:safety} (\emph{Crosslink Guide}, \S``Coverage
of the formal models'').

\begin{table}[ht]
\caption{Status by goal (\emph{Crosslink Guide}, \S``Status by goal and by
rule'').}
\label{tab:status}
\begin{tabular}{@{}p{0.36\textwidth}p{0.36\textwidth}p{0.22\textwidth}@{}}
\toprule
Goal & Result & Status \\
\midrule
$\fin_i$ never rolls back & Lemma~\ref{lem:fin-linear} & proved (Crosslink
  Guide) \\
$\fin_i^t \preceqbc (\ba_\mu)_i^t$ & Proposition~\ref{prop:ledger-prefix} &
  proved (TFL book) \\
Assured Finality & Theorem~\ref{thm:pa-af}, Theorem~\ref{thm:fa-af},
  Corollary~\ref{cor:af-either} & proved (TFL book); proved (Crosslink
  Guide) \\
Safety with the BFT protocol subverted & Corollary~\ref{cor:bft-subversion} &
  proved (Crosslink Guide) \\
Bounded availability & Proposition~\ref{prop:bounded-run}, finalised chain
  on the best chain & proved (Crosslink Guide); otherwise open problem \\
Liveness & \S\ref{sec:live} & argued (TFL book) \\
The premises for Zcash & Remark~\ref{rem:pow-prefix}; no BFT protocol fixed
  & open problem \\
\bottomrule
\end{tabular}
\end{table}

\end{document}
